\subsection{Removing the central ambiguity: the full linear modular action}

\begin{claim}{C.a15}{Proved}{The unique cusp-preserving lift from the projective modular group}
Write
\[
 H=\widehat{\mathrm{SL}_2(\mathbb Z)}
  =\langle S,R\mid S^4=1,\ S^2=R^3\rangle_{\rm prof},\quad
 z=S^2=R^3=-I,
\]
$P=H/\langle z\rangle=\langle s,r\mid s^2=r^3=1\rangle_{\rm prof}$,
and $T=SR$, $e=sr$. Let $u\in\operatorname{Aut}(P)$ satisfy
\[
 u(e)=f^{-1}e^\mu f,\qquad\mu\in\widehat{\mathbb Z}^{\times},\quad f\in P.
 \tag{C.a15.1}
\]
For the linear congruence coordinates we choose
\[
 q_H(S)=-S_0=\begin{pmatrix}0&-1\\1&0\end{pmatrix},\qquad
 q_H(R)=\begin{pmatrix}0&1\\-1&1\end{pmatrix},\qquad
 q_H(T)=\begin{pmatrix}1&-1\\0&1\end{pmatrix}=E,
\]
where $S_0$ is the matrix used for the projective involution in C.d1.
Changing its lift from $S_0$ to $-S_0$ leaves every projective
cocycle and conjugation formula of C.d1 unchanged and makes the
chosen full cusp generator map to the unipotent $E$.
Then precisely one of the two lifts of $u$ to $\operatorname{Aut}(H)$
satisfies
\[
 \widetilde u(T)=\widetilde f^{-1}T^\mu\widetilde f,
 \tag{C.a15.2}
\]
where $\widetilde f$ is any lift of $f$. It sends $S$ to a conjugate
of $S^{\mu\bmod4}$ and $R$ to a conjugate of $R^{\mu\bmod6}$.
The two lifts differ by the central transvection of C.a7. The lift
(C.a15.2) respects composition and inner automorphisms, and hence induces
an action of $\widehat{\mathrm{GT}}_0$ on $H$ by outer automorphisms.
\end{claim}
% anchor:C.a15
\manuscriptanchor{\S 36; fols.266--270 verso; cote140-2 scans521--530; \S \S 45--46; fols.387--391; cote140-3 scans176--184}

\begin{proof}
Use the same finite-torsion conjugacy theorem as in C.d1.
The element $u(s)$ is conjugate to $s$, and $u(r)$ to either $r$ or
$r^{-1}$. Abelianizing (C.a15.1) in $P^{\rm ab}=C_2\times C_3$ shows
which: $u(r)$ is conjugate to $r^{\mu\bmod3}$. Thus choose $a,b\in P$
with $u(s)=asa^{-1}$ and $u(r)=br^{\mu}b^{-1}$, and choose lifts
$\widetilde a,\widetilde b\in H$. The assignments
\[
 S\longmapsto\widetilde a S^\mu\widetilde a^{-1},\qquad
 R\longmapsto\widetilde b R^\mu\widetilde b^{-1}
 \tag{C.a15.3}
\]
respect the displayed presentation of $H$: the square of the first
image and cube of the second are both $z$, since $\mu$ is odd. They
therefore define a continuous endomorphism $\widetilde u$ with quotient
$u$ and with $z$ fixed. It is an automorphism: its kernel maps trivially
to $P$ and meets $\langle z\rangle$ trivially; for surjectivity lift a
preimage in $P$ and correct by $z$ in the center. Compactness makes the
resulting bijection a topological isomorphism.

The quotient equation implies
$\widetilde u(T)=z^\epsilon\widetilde f^{-1}T^\mu\widetilde f$
for some $\epsilon\in\{0,1\}$. In the cyclic group
$H^{\rm ab}=\mathbb Z/12$, the class $t$ of $T$ generates, the class
of $S$ is $3t$, that of $R$ is $-2t$, and that of $z$ is $6t$.
These identities follow by abelianizing $2S=3R$ and $4S=0$; their
converse defines a surjection onto $\mathbb Z/12$, so no extra relation
has been imposed. Formula (C.a15.3) therefore sends $t$ to
$3\mu t-2\mu t=\mu t$. Comparison with the previous display gives
$6\epsilon=0\pmod {12}$, so $\epsilon=0$. In particular the result is
independent of the choices of $a,b$ and their lifts.

Any two lifts of $u$ differ by a character $H\to C_2$ fixing $z$,
as in C.a7. There are exactly two; the nontrivial one is $\eta$ and
$\eta(T)=1$. It changes the right side of (C.a15.2) by $z$, so only one
lift satisfies that equation. The lift of a composite satisfies the same
condition with multiplier the product of the two multipliers; uniqueness
proves compatibility with composition. A lift of an inner automorphism
of $P$ is the inner automorphism defined by a lift of its conjugator,
and satisfies the cusp condition; therefore the construction descends
to outer classes. Apply it to the automorphisms of $P$ in B.m1--B.m2.
\end{proof}
The uniqueness in this assertion is a uniqueness of a central lift with
its full cusp generator specified. It does not identify or split the
different, noncentral extension asked about in Part D. On outer classes
one can equivalently require that $T$ be sent to a conjugate of $T^\mu$:
$T^\mu$ and $zT^\mu$ have different abelianizations, so they cannot
be conjugate.

\begin{claim}{C.a16}{Proved}{The congruence descent holds in the full linear group}
Let $q_H:H\twoheadrightarrow\mathrm{SL}_2(\widehat{\mathbb Z})$ be
the full congruence homomorphism, and take the matrices $U(u)$ of C.d1.
For the cusp-preserving lift of C.a15,
\[
 q_H\circ\widetilde u=\operatorname{Int}(U(u))\circ q_H.
 \tag{C.a16.1}
\]
Thus the full, not only projective, congruence kernels are preserved by
this action of $\widehat{\mathrm{GT}}_0$. Its outer congruence image,
that of $\widehat{\mathrm{GT}}$, and the arithmetic image agree at
every finite level: they are the conjugations by
$\mathrm{GL}_2(\mathbb Z/n)$ modulo the inner conjugations by
$\mathrm{SL}_2(\mathbb Z/n)$, with determinant parameter
$(\mathbb Z/n)^{\times}/((\mathbb Z/n)^{\times})^2$.
\end{claim}
% anchor:C.a16
\manuscriptanchor{\S 39; fols.298 verso--301; cote140-3 scans21--26; \S 41; fols.315--316 verso; scans54--57}

\begin{proof}
First normalize $u_0$ so that $u_0(e)=e^\mu$. In C.d1 its matrix is
\[
 V=\begin{pmatrix}\mu&\mu B\\0&1\end{pmatrix},\qquad
 q_H(T)=E=\begin{pmatrix}1&-1\\0&1\end{pmatrix}.
\]
The two continuous homomorphisms $q_H\circ\widetilde u_0$ and
$\operatorname{Int}(V)\circ q_H$ have equal images after passage to
the quotient by the single global $\{\pm I\}$, by C.d1. Their
pointwise ratio is therefore a continuous character
$\kappa:H\to\{\pm I\}$: centrality makes this ratio multiplicative.
But C.a15 gives $q_H(\widetilde u_0(T))=E^\mu$, while direct matrix
multiplication gives $VEV^{-1}=E^\mu$. Hence $\kappa(T)=1$.
Every character to $C_2$ factors through $H^{\rm ab}=C_{12}$, generated
by $T$, and consequently $\kappa=1$. This proves (C.a16.1) for $u_0$.
Removing the normalization by the same inner conjugator as in C.d1
proves it for $u$. The alternative lift differs by $\theta$ and does
not satisfy this equality: it sends $T$ to $zT^\mu$.

At level $n$, equality (C.a16.1) proves descent before an action on the
quotient is asserted. A matrix commuting with all of
$\mathrm{SL}_2(\mathbb Z/n)$ commutes in particular with both elementary
unipotents; comparing entries shows that it is a scalar matrix.
Thus two $\mathrm{GL}_2$ conjugations differ by an inner $\mathrm{SL}_2$
conjugation exactly when their determinants differ by the square of a
unit. Reduction $\mathrm{SL}_2(\mathbb Z)\to
\mathrm{SL}_2(\mathbb Z/n)$ is surjective, as follows from elementary
row operations over each prime-power local factor and the Chinese
remainder theorem; hence these are exactly the inner conjugations
coming from $H$ as well. C.d1 gives $\det U(u)=\mu(u)$.
The cyclotomic character $G_{\mathbb Q}\to\widehat{\mathbb Z}^{\times}$
is surjective, so its finite residue characters realize every determinant
square class. The inclusions of the arithmetic image into
$\widehat{\mathrm{GT}}\subseteq\widehat{\mathrm{GT}}_0$ then force
equality of all three finite outer images. The arithmetic lift is the
cusp-preserving one because geometric tame inertia at the modular cusp
is acted on by the cyclotomic character; the Tate family fixes this
full inertia subgroup over $\mathbb Q$. Equivalently its homological
congruence action sends a unipotent to a conjugate unipotent, which
excludes the alternative $zT^\mu$ already over characteristic zero.
\end{proof}
